\section{OpenAVT Design}
\label{sec:design}
\subsection{Architecture and separation of authority}
\sys{} separates a protected control exchange from ordinary conversational
speech (Figure~\ref{fig:architecture}). The public \code{VoiceSession} API owns
the audio adapter, authentication lifecycle, and serialized request exchange.
A receiving \code{Gateway} owns an allowlist of handlers, request validators,
and an authorization policy. The same gateway operates over both transport
profiles. The model has no entry point that directly invokes a registered
business handler while bypassing this gateway.

The prototype's Agent Voice Authority (AVA) performs account registration,
CSR enrollment, key-possession proof, certificate issuance, renewal, and
revocation. A distinct business signer issues customer grants. Separating these
keys prevents possession of an enrollment CA key from being treated as the
intended application permission. A production service would additionally need
credible account authentication, key recovery, issuer governance, and appropriate
identity assurance. The local authority is not a production CA.

\begin{figure*}[t]
\centering
\begin{tikzpicture}[font=\small,>=Latex,
 box/.style={draw,rounded corners,align=center,minimum height=10mm,text width=29mm},
 flow/.style={->,thick}]
\node[box] (model) at (0,0) {Agent model\\proposes intent};
\node[box] (sdk) at (3.6,0) {Session + key store\\checks channel};
\node[box] (modema) at (7.2,0) {TLS + ARQ or QUIC\\audio modem};
\node[box] (modemb) at (10.8,0) {Audio modem\\protected channel};
\node[box] (gateway) at (14.4,0) {Gateway + policy\\business handler};
\draw[flow] (model) -- (sdk);
\draw[flow] (sdk) -- (modema);
\draw[flow,<->] (modema) -- node[above=7mm,align=center]{call audio} (modemb);
\draw[flow] (modemb) -- (gateway);
\node[box] (grant) at (5.4,2.0) {Business grant signer\\user permission};
\draw[flow] (grant) -- node[left]{grant} (sdk);
\draw[flow,dashed] (grant.east) -| node[pos=0.3,above]{verification / current state} (gateway.north);
\node[draw,dashed,align=center,text width=28mm] (speech) at (10.8,-1.6) {Ordinary speech\\remains untrusted};
\draw[dashed,->] (speech) -- (modemb);
\end{tikzpicture}
\caption{Protected structured control travels through the call's actual audio. Enrollment and grant acquisition use separate services. The gateway checks permission outside the model; the inner channel does not authenticate the semantic content of ordinary speech. Dashed edges denote separate trust or service interactions.}
\label{fig:architecture}
\end{figure*}

\subsection{Audio modulation and framing}
The \code{voiceband} profile uses binary FSK at 1 kbps, with 1 and 2 kHz tones
in 8 kHz mono samples. Its streaming decoder acquires a preamble at unknown
sample offsets, reads a length-bounded frame, and verifies a CRC. It assumes
stable timing within each frame and has no clock-drift recovery.

The \code{pcm} profile maps seven bits to one of 128 selected G.711 mu-law
reconstruction levels per sample, giving a raw signaling rate of
$8000\times7=56{,}000$ bit/s on a sample-preserving digital path. Preamble,
header, CRC, padding, and transport cryptography reduce application throughput.
This is not a claim of 56 kbps through an arbitrary voice codec, acoustic link,
or PSTN route. Transcoding and voice enhancement can alter those exact levels.

Both profiles use the repository's sequence, frame-kind, and length header.
CRC detects accidental damage but supplies no cryptographic authenticity.
The cryptographic layer must reject modified bytes even when an attacker
recomputes a valid CRC. Call adapters deliver 8 kHz mono integer samples;
WebRTC uses audio RTP rather than a data channel, and the provider adapter
wraps G.711 audio in the provider's media envelope.

\subsection{TLS over reliable audio}
The original channel uses OpenSSL \code{MemoryBIO} endpoints configured for
mutual TLS 1.3~\cite{rfc8446}. It fragments outgoing TLS bytes into bounded
frames, transmits a frame, waits for its audio acknowledgment, and retries on
timeout. Sequence numbers suppress duplicate delivery. Lost frames or ACKs
therefore do not require a lossless physical medium.

This profile is a useful baseline: it is small, uses a mature TLS implementation,
and offers transparent ordered bytes. Its stop-and-wait operation can leave
capacity idle while a low-rate acknowledgment is played. Our evaluation must
measure that cost against QUIC's larger handshake, packet overhead, and recovery
behavior instead of assuming an architectural win.

\subsection{QUIC over unreliable audio datagrams}
The QUIC channel uses aioquic 1.3.0's sans-I/O packet engine~\cite{aioquic}.
It supplies received packets and a monotonic clock, drives the engine's timers,
and modulates generated packets. The routing address passed to the engine is
an opaque local token; the channel creates no UDP socket. This uses QUIC's
packet format, TLS handshake, reliable stream machinery, and recovery
algorithms~\cite{rfc9000,rfc9001,rfc9002} over an experimental non-UDP underlay.
It is not standard Internet QUIC interoperability with an unmodified endpoint.

Packets are limited to 1200 bytes. Audio frame kind 3 carries a fragment index,
fragment count, total packet length, and packet bytes. The four-byte fragment
header leaves 252 bytes per FSK frame, requiring up to five fragments; the PCM
profile fits a packet into one frame. Each packet has its own audio sequence.
The receiver accepts out-of-order fragments and duplicates, bounds reassembly
to 16 partial packets, and expires incomplete entries after 5 seconds for PCM
or 60 seconds for FSK. It discards missing or invalid fragments rather than
requesting audio retransmission. A later packet does not wait for an earlier
sequence gap. These are link datagrams carrying QUIC packets, not use of QUIC's
optional unreliable application DATAGRAM extension.

The engine and audio sender use separate asynchronous tasks. The engine can
process inbound packets while a bounded output burst is playing; it builds
further output and services timers between bursts. Initial RTT is 0.5 seconds
for PCM and 25 seconds for FSK to accommodate audio serialization. These values
are prototype defaults, not optimized parameters. Native congestion control
limits generated bursts; a 128-entry output queue supplies an additional bound.
The application serializes requests on bidirectional stream 0. We do not yet
exploit multiplexing and make no claim about avoiding application head-of-line
blocking across concurrent operations.

\subsection{Mutual authentication and session lifecycle}
Both profiles require an explicitly selected root and a client-specified
service hostname. QUIC uses version 1, ALPN \code{openavt/1}, and TLS 1.3 with
no configured resumption tickets or 0-RTT. The QUIC integration additionally
checks both leaf certificates' chain, validity, expected purpose,
digital-signature usage, and non-CA status. TLS CertificateVerify supplies
the proof of key possession.

Requiring a client certificate needs care: aioquic's default server setup does
not expose a public required-client-certificate option. A narrow subclass
enables its certificate-request hook and retrieves the peer certificate for
explicit validation before the application is marked verified. The client
also waits for a fixed encrypted server confirmation on unidirectional stream
3, sent only after the server's client validation. This prevents the client
from releasing its grants merely because it has authenticated the server.
The subclass also reads stream acknowledgment offsets and pending ACK state
to bound delivery waits. These private interfaces are a maintenance and audit
liability; the dependency is pinned and upgrades need the negative handshake
and teardown tests. This is not a new key-exchange protocol or a formal proof
of the library integration.

A session moves from new to authenticating to verified, or fails and closes.
Certificate or transport failure blocks protected operations and invalidates
the risk monitor's authenticated context. Application JSON messages are limited
to 64 KiB, and message queues to 32 entries. Protected logs omit application
contents and secrets; QUIC telemetry retains loss counters rather than
unbounded packet traces or traffic keys.

The shared audio scheduler reserves modem bursts and prevents speech mixing.
Speech resumes between protected exchanges, and the receiver routes samples
outside detected modem activity to the optional speech adapter. This heuristic
is useful scheduling, not a proof of semantic isolation: incomplete or
unrecognized modem activity can still require careful handling by an ASR stack.

\subsection{Delegation and execution}
The client requests a live challenge before submitting an operation. The
gateway binds the response to a request identifier and nonce, validates an
allowlisted action and its schema, and invokes its authorization policy.
The standalone AVA grant signs issuer, customer, audience, agent identity,
scope, resource, validity interval, unique grant ID, and certificate fingerprint.
Its fixed read-only policy rejects additional arguments. The challenge and
grant are checked and consumed in persistent transaction state. A new action
does not inherit permission from a successful earlier read.

The optional PAP bridge obtains a grant over HTTPS using a current session and
DPoP-style key possession. It also proves possession of the voice key and binds
the grant to the live voice challenge, request ID, action, resource, customer,
and canonical argument hash. Its lifetime is at most 30 seconds and no longer
than the session or certificate. The gateway verifies the current permission
and consumes replay state. This policy supports a second permission check
after risk assessment, immediately before the handler. The standalone AVA
policy does not currently implement that second hook; production policies
must define their revocation and commit-time race semantics explicitly.

The risk adapter receives bounded audio observation and structured request
signals. A block or unavailable policy can reduce execution; an ``allow''
assessment cannot bypass the grant. Policies, schemas, and business state
remain outside the language model. Keys and account tokens are kept out of
model context, although applications must separately control any confidential
business response later supplied to a model or spoken aloud.

Finally, transport deduplication is not application exactly-once execution.
If a submitted request loses its confirmed response, the SDK raises
\code{RequestOutcomeUnknown} and closes the call. The PAP demonstration proposes
a write over audio but confirms a stable operation ID through HTTPS; repeated
confirmation and status recovery produce one simulated effect. Other writes
need equivalent transaction or idempotency semantics.

\subsection{Teardown as part of delivery}
A receiver can obtain the final response before its delayed QUIC ACK has been
generated or finished audio playback. Closing at that point can leave the
sender uncertain about delivery. The integration waits for outgoing stream
bytes to be acknowledged, for pending ACK generation, and for queued audio
to finish. A regression test holds the final ACK waveform and checks that
drain waits. This bug was exposed when validating the slower FSK path; it
illustrates how ordinary network-lifecycle assumptions can fail under audio
serialization.
